To accurately model the complex socio-ecological dynamics of municipal waste management, the simulation operates on a highly structured dual timescale and a spatially explicit environment \citep{Brugiere2022}.

\paragraph{Spatially Explicit Grid Architecture}
The physical environment is represented as a $50 \times 50$ \texttt{MultiGrid}. Crucially, the grid is configured as **non-toroidal**, ensuring that agents cannot "wrap around" the edges of the municipality—a constraint that preserves the geographical integrity of the seven distinct barangays. This spatial layer facilitates the explicit movement of \texttt{EnforcementAgent} objects and calculates the local social norm influence based on the density of compliant households within a defined Moore neighborhood \citep{TianReview2024}.

\paragraph{Dual-Clock Temporal Resolution}
The simulation utilizes a dual-timescale architecture to capture the friction between micro-level human behavior and macro-level administrative governance \citep{Akbarpour2021}:
\begin{itemize}
    \item \textbf{Daily Micro-Ticks:} At the micro-level, the environment processes daily iterations (1 tick = 1 day). During each tick, \texttt{HouseholdAgents} execute routine interactions, evaluate behavioral friction via their TPB utility functions, and face the stochastic probability of being caught by patrolling \texttt{EnforcementAgents}.
    \item \textbf{Quarterly Policy Epochs:} The macro-level operates on a 90-day epoch. Every 90 ticks, the simulation execution pauses to allow the \texttt{MayorAgent} to interface with the Deep Reinforcement Learning (DRL) policy.
\end{itemize}
The total simulation horizon is fixed at $T = 1,080$ ticks (50,000 during training), representing the standard three-year political term of Philippine local officials.



\subsection{Polycentric Governance: The Coupled ABM-RL Loop}

The model implements a polycentric governance structure as mandated by R.A. 7160 and R.A. 9003 \citep{RA9003, Ostrom2010}. The interaction between these layers is governed by a closed-loop Markov Decision Process (MDP) where the \texttt{MayorAgent} acts as the policy actuator.

\paragraph{The HuDRL Action Space}
Every 90 days, the Mayor Agent observes a 17-dimensional state vector and outputs a 21-dimensional action vector $\mathcal{A} \in [-1, 1]^{21}$ \citep{TianReview2024}. This vector is partitioned into triplets for each of the seven barangays:
\begin{equation}
    \mathcal{A}_{B_i} = \{ \text{Enforcement Allocation, Incentive Allocation, IEC Allocation} \}
\end{equation}
\addequation{The Action Space Equation}
This allows the AI to develop highly localized "triage" strategies—for instance, deploying heavy enforcement in low-compliance areas like \textit{Poblacion} while shifting to education-heavy strategies in high-compliance areas like \textit{Binuni}.

\paragraph{Reward Shaping and Redemption Reset}
To prevent "policy inertia," the simulation enforces a quarterly \textit{Redemption Status Reset}. Every 90 days, all households' incentive eligibility flags are cleared. Households must re-demonstrate compliance within the new quarter to qualify for municipal rewards. 

Computationally, the HuDRL's performance is evaluated via a complex reward function $R_t$ defined in the \texttt{BacolodGymEnv}. The reward function prioritizes "Jackpot Rewards" for compliance gains while penalizing "Administrative Waste" (over-funding saturated areas) and "Political Collapse" (when approval drops below 10\%).

\vspace{0.5cm} 
\begin{algorithm}[H]
\caption{The Coupled ABM-RL Main Simulation Loop (Plain English)}
\label{alg:main_loop_english}
\begin{algorithmic}[1]
\State \textbf{Setup:} Create a 50 by 50 flat map and set the starting yearly budget to \textpeso 1.5 million.
\State \textbf{Instantiation:} Create 7 Barangay agents and populate the map with household agents using real census numbers.
\For{every single day from Day 1 to Day 1080 (roughly 3 years)}
    \Statex \textbf{// Phase 1: Quarterly AI Planning}
    \If{it is the end of the quarter (every 90 days)}
        \State \parbox[t]{0.88\linewidth}{Gather the current situation: check people's compliance, attitudes, remaining budget, and the Mayor's political support.\strut}
        \State \parbox[t]{0.88\linewidth}{Ask the AI to decide the next best action based on this current situation.\strut}
        \State \parbox[t]{0.88\linewidth}{Convert the AI's decision into a new budget plan for the 7 Barangays.\strut}
        \State \parbox[t]{0.88\linewidth}{\textbf{Reset:} Clear all household reward trackers to start fresh for the new quarter.\strut}
    \EndIf
    \Statex \textbf{// Phase 2: Daily City Life}
    \For{each Barangay}
        \State \parbox[t]{0.88\linewidth}{Carry out daily tasks like sending out local patrols (Tanods) and making public announcements (Recorida).\strut}
    \EndFor
    \For{each Household}
        \State \parbox[t]{0.88\linewidth}{Weigh the decision to segregate trash by comparing community peer pressure against the fear of getting fined.\strut}
        \State \parbox[t]{0.88\linewidth}{Decide to either follow the rules (Compliant) or ignore them (Non-compliant).\strut}
    \EndFor
    \Statex \textbf{// Phase 3: Enforcement and Consequences}
    \State \parbox[t]{0.88\linewidth}{Enforcers patrol the map, catch violators, and collect fines (which get added back to the budget).\strut}
    \State \parbox[t]{0.88\linewidth}{Update the Mayor's political support (fining people too harshly lowers support).\strut}
    \If{the Mayor's political support drops to 10\% or lower} 
        \State \textbf{Stop the simulation:} The local government has collapsed due to low public support.
    \EndIf
\EndFor
\end{algorithmic}
\end{algorithm}